\section{Limitations}
\label{sec:limitations}

\textbf{Measurements.} The sequential times in Table~\ref{tab:sequential} come from one
machine, with \SeqRunsPhrase. \SeqRepeatSentence{} Small differences between neighbouring
rows therefore carry no meaning. The conclusions rest on ratios far larger than that.
\MpiNoiseSentence

\textbf{The tables and the processor's caches.} The two pair tables are 3.5~MB each. Whether
they fit in a core's cache decides how long a lookup takes, above all in the matrix product,
where the additions wait for each other. The speedup therefore depends on the processor, and
the figures here should not be read as what another machine will show.

\textbf{The sequential reference in the MPI comparison.} The MPI program times the
assignment's elimination on an array whose size is chosen at run time, so that one binary
can sweep all sizes. The assignment program fixes $N$ at compile time and checks every
index. \MpiReferenceSentence

\textbf{Where the MPI processes ran.} \MpiNodesLimitation

\textbf{Scope of the cache.} The tables need about $4q^2$ bytes, which is why fields above 4096
elements are forwarded to the original arithmetic. Constructing elements is not cached. The
matrix product is still sequential and is now the largest part of the program's run time.
